\section{Draw the allowed input window}\label{packet-x-window}
\begin{figure}[htbp]
\centering
\begin{tikzpicture}[alt={A row of twelve byte cells shows offsets zero through eleven. The first eleven are inside an available interval of eleven bytes; the checksum cell at offset eleven lies outside it. A four-byte read starting at offset eight spans payload bytes eight through ten and the unavailable checksum, so its entire width does not fit.},x=10mm,y=10mm]
\foreach \i in {0,...,11} {
 \draw (\i,0) rectangle +(1,1);
 \node[font=\scriptsize] at (\i+0.5,0.5) {\i};
}
\draw[very thick] (0,1.25)--(11,1.25);
\draw[very thick] (11,-0.2)--(11,1.5);
\node[above] at (5.5,1.3) {Eleven available bytes};
\draw[<->] (8,-0.4)--(12,-0.4);
\node[below] at (10,-0.5) {Four-byte read};
\node[above,anchor=east] at (12,2) {End boundary at offset 11};
\end{tikzpicture}
\caption{A load can begin inside the input and still cross its end. Cell numbers are offsets, not byte values.}
\end{figure}
The default packet needs twelve bytes. With only eleven available, offsets zero
through ten may be read, but the trailer at offset eleven may not. A four-byte
load from offset eight would overlap that forbidden byte. The preceding availability
check must return need more before this implementation reaches its checksum scan.
The diagram is native explanatory artwork, separate from the pending IDE captures.
